\documentclass[11pt]{article}
\usepackage[a4paper,margin=27mm]{geometry}
\usepackage{amsmath,amssymb,amsthm,hyperref}
\newtheorem{lemma}{Lemma}\newtheorem{proposition}{Proposition}
\newcommand{\I}{\mathcal I}\newcommand{\E}{\mathcal E}\newcommand{\C}{\mathcal C}
\title{A central Frobenius--Schur obstruction for the relative center of $R6_{8,1660}$}
\author{}\date{30 September 2026}
\begin{document}\maketitle
\begin{lemma}[A central factor and the second indicator]
Let $\mathcal A$ be a pivotal fusion category, let $b$ be a self-dual
simple object of $Z(\mathcal A)$ whose underlying object is simple, and
let $t$ be self-dual and simple in $\mathcal A$. If $bt$ is simple,
then $\nu_2(bt)=\nu_2(b)\nu_2(t)$.
\end{lemma}
\begin{proof}
We suppress coherent associators and duality constraints. Our
half-braiding convention is $c_{b,x}:bx\longrightarrow xb$.
Choose self-dualities $f_b:b\to b^*$ in $Z(\mathcal A)$ and
$f_t:t\to t^*$ in $\mathcal A$. With $j$ denoting the pivotal structure,
these satisfy $f_b^*j_b=\nu_2(b)f_b$ and
$f_t^*j_t=\nu_2(t)f_t$. A self-duality for $bt$ is
\[
 f=c_{b^*,t^*}(f_b\otimes f_t):bt\longrightarrow t^*b^*.
\]
Duality of half-braidings gives
$(c_{b^*,t^*})^*=c_{b^{**},t^{**}}$. Hence
\begin{align*}
 f^*j_{bt}
 &=(f_t^*\otimes f_b^*)c_{b^{**},t^{**}}(j_b\otimes j_t)\\
 &=(f_t^*j_t\otimes f_b^*j_b)c_{b,t}\\
 &=\nu_2(b)\nu_2(t)(f_t\otimes f_b)c_{b,t}\\
 &=\nu_2(b)\nu_2(t)c_{b^*,t^*}(f_b\otimes f_t).
\end{align*}
The second equality uses pivotal naturality, and the last uses naturality
in $t$ and the fact that $f_b$ is a morphism of central objects. Since
$bt$ is simple, this is the defining second-indicator scalar.
No half-braiding for $t$ is assumed.
\end{proof}

Use the established pseudo-unitary reduction for $R6_{8,1660}$ and
its indicators
\[
 \nu_2(g)=\nu_2(X)=\nu_2(gX)=\nu_2(\sigma X)=1,
 \qquad\nu_2(\sigma)=-1.
\]
The relative center $\E=Z_{\I}(\C)$ has the previously proved ring
$B[t]/(t^2-L-At)$, where $B=K_0(Z(\I))$, $L=1$ or
$e=(g,g)$, and $t^*=Lt$. The reverse-braided central lift
$Z(\I)^{\rm rev}\to Z(\E)$ is fully faithful. Give $\E$ the induced
positive spherical structure.

\begin{proposition}
The branch $L=1$ is impossible. Thus precisely the following two
previously listed ring representatives remain as necessary possibilities:
\begin{align*}
 L&=e,& A&=\sum_{a,b\in\{1,g,\sigma\}}(a,b),\tag{I}\\
 L&=e,& A&=2(\sigma,1)+2(\sigma,g)+3(\sigma,\sigma).\tag{II}
\end{align*}
Interchanging the Ising factors gives the third labelled alternative.
\end{proposition}
\begin{proof}
Suppose $L=1$. Then $t$ and $bt$, with $b=(\sigma,1)$, are self-dual
simples. Under the pivotal forgetful functor $F:\E\to\C$, they map to
$X$ and $\sigma X$, respectively, while $b$ maps to $\sigma$.
For each of these three simples, the source and target spaces
$\operatorname{Hom}(Y,Y^*)$ are one-dimensional; the faithful pivotal
functor identifies their transpose operators. Their indicators are
therefore $1,1,-1$, respectively. The lemma instead gives
$\nu_2(bt)=-1$, a contradiction.

Of the $58$ labelled rings in the established arithmetic reduction,
$55$ have $L=1$. The remaining three are the two displayed representatives
and the factor exchange of (II). This extraction is recorded, with
its input hash, in \texttt{structural\_cuts\_certificate.json}.
\end{proof}

For either surviving branch, among the nine complementary simples
$bt$ only $S=(\sigma,\sigma)t$ is self-dual; all other eight form four
dual pairs. Furthermore $\nu_2(S)=1$. Indeed,
$F(S)=X\oplus gX$ has two distinct self-dual constituents of indicator
$1$, each with multiplicity one. A self-duality of $S$ maps to a
nondegenerate self-duality on that direct sum with the same transpose
sign. Its two nonzero diagonal blocks force that sign to be $1$.
This is not an assertion that arbitrary tensor functors preserve
Frobenius--Schur indicators.

\begin{proposition}[Excluding the all-ones coefficient]
Only representative (II), up to exchanging the Ising factors, remains.
\end{proposition}
\begin{proof}
The relative center $\E$ is Morita equivalent to
$\C\boxtimes\I^{\rm op}$: regard $\C$ as a module category under
left multiplication by $\C$ and right multiplication by $\I$.
An endofunctor commuting with the left action is right tensoring by
an object of $\C$, and its compatibility with the right action is
exactly a relative half-braiding. Consequently
\[
 |\operatorname{Inv}Z(\E)|=4|\operatorname{Inv}Z(\C)|.
\]
All four invertibles of $\E$ belong to $B$ and have canonical central
lifts. Therefore the forgetful homomorphism
$\operatorname{Inv}Z(\E)\to\operatorname{Inv}(\E)$ is surjective;
its kernel is the group of tensor automorphisms of the identity,
namely the character group of the universal grading group $U(\E)$.

In case (I) every basis object has trivial universal degree. Indeed,
$t^2$ contains $t$, forcing $(\deg t)^2=\deg t$ and hence $\deg t=1$; since $At$
contains every complementary simple, all $b\in B$ have degree $1$.
Thus $|\operatorname{Inv}Z(\E)|=4$, which forces
$|\operatorname{Inv}Z(\C)|=1$.

But the previously established exhaustive restriction-profile domain
always contains exactly one central simple restricting to $g$ alone.
This is the high vector $(1,0,0,0,0)$, where the omitted first
coordinate is the unit multiplicity zero. The multiplicity is one
in every raw profile in both low frames, even before later filters:
$371$ profiles for each of two labellings in frame zero and $2011$
for each of two labellings in frame one. Its dimension is one,
so it is a nontrivial invertible of $Z(\C)$, a contradiction.
The tiny exact checker \texttt{verify\_structural\_cuts.py}
checks all $4764$ raw profiles and records their hashes and indices.
\end{proof}

In case (II), put $p=(g,1)$, $q=(1,g)$, and $Y=(\sigma,1)t$.
The universal grading is $\mathbb Z/2$: the first Ising spin and
$t$ have odd degree, while the second Ising factor has trivial degree.
Its neutral component has nine simples and is generated by
$p,q,\sigma_2,Y$ with
\[
 p^2=q^2=1,\quad q\sigma_2=\sigma_2,
 \quad\sigma_2^2=1+q,\quad pY=Y,\quad Y^*=qY,
\]
\[
 Y^2=q(1+p)+\bigl(4(1+q)+6\sigma_2\bigr)Y.
\]
The second Ising factor has its canonical central lift.
This remaining ring is a necessary possibility only. No existence
or exclusion of the original rank-six ring is asserted here.
\end{document}
